% exam.tex — Documento-mãe do PDF de simulado offline.
% \documentclass[11pt,a4paper]{article}
\usepackage{fontspec}
\setmainfont{PlusJakartaSans-Regular.ttf}[
  Path = /opt/fonts/PlusJakartaSans/,
  BoldFont = PlusJakartaSans-Bold.ttf
]
\newfontfamily\jakartamedium{PlusJakartaSans-Medium.ttf}[Path = /opt/fonts/PlusJakartaSans/]
\usepackage{polyglossia}
\setmainlanguage{portuguese}
\usepackage{ucharclasses}
\newfontfamily\greekfallback{NotoSans-Regular.ttf}[Path = /opt/fonts/NotoSans/]
\setTransitionsForGreek{\begingroup\greekfallback}{\endgroup}
\usepackage{microtype}
\usepackage{graphicx}
\usepackage[margin=2cm]{geometry}
\usepackage{fancyhdr}
\usepackage{xcolor}
\usepackage{enumitem}
\usepackage{needspace}
\usepackage[most]{tcolorbox}
\usepackage{tikz}

% Cores exatas equivalentes ao motor pdf-lib atual (supabase/functions/generate-exam-pdf/index.ts:29-36)
\definecolor{winedark}{HTML}{591230}
\definecolor{winemid}{HTML}{6B1C3D}
\definecolor{winelight}{HTML}{8C2E4C}
\definecolor{graylight}{HTML}{EBEBEB}
\definecolor{graymid}{HTML}{8C8C8C}
\definecolor{graybg}{HTML}{F5F5F5}
\definecolor{titlesubtitle}{HTML}{FFD9E6}
\definecolor{titlesubtitle2}{HTML}{FFB2C7}
\definecolor{rulesboxtint}{HTML}{FAF0F2}

\pagestyle{fancy}
\fancyhf{}
% cabeçalho/rodapé completos entram junto com exam.tex (Task 5) — aqui só reset de estilo

% Exceções de hifenização para termos médicos compostos que o pacote de
% PT-BR erra sozinho. Um item por entrada, sílabas separadas por hífen nos
% pontos de quebra válidos. Vazio por padrão — popular sob demanda conforme
% casos ruins aparecerem em revisão editorial.
% Exemplo (comentado): \hyphenation{ele-tro-car-di-o-gra-ma}
\hyphenation{}
 (Task 3) define classe, fontes, cores, pacotes.
% % cover.tex — Capa (recriação visual de drawCoverPage, supabase/functions/generate-exam-pdf/legacyPdfLib.ts:200-414)
% Inputted from exam.tex (Task 5) logo após \begin{document}. Placeholders @@...@@
% são substituídos por dados reais na Task 9 (Node.js). NÃO editar valores aqui —
% este arquivo permanece com os placeholders literais.

% ── Bloco de título com fundo wine ──
\begin{tcolorbox}[colback=winedark, colframe=winemid, sharp corners, boxrule=1pt, width=\textwidth, left=8pt, right=8pt, top=10pt, bottom=10pt]
  {\fontsize{26}{30}\selectfont\bfseries\color{white} PROVA OFFLINE}\\[4pt]
  {\fontsize{14}{16}\selectfont\color{titlesubtitle} @@SIMULADO_TITLE@@}\\[2pt]
  {\fontsize{10}{12}\selectfont\color{titlesubtitle2} Modo Offline · Gabarito Digital}
\end{tcolorbox}

\vspace{18pt}

% ── 3 info-cards lado a lado: Questões / Duração / Alternativas ──
% Nota: \tcbitem só é definido pelo ambiente tcbitemize (que envolve tcbraster
% internamente) — não por tcbraster diretamente. Mesmas opções do brief, apenas
% o nome do ambiente corrigido.
\begin{tcbitemize}[raster columns=3, raster equal height, colback=graybg, colframe=graylight, rounded corners, boxrule=0.5pt]
  \tcbitem {\small\color{graymid} Questões}\\ {\Large\bfseries\color{winedark} @@QUESTIONS_COUNT@@}\\ {\tiny\color{graymid} Múltipla escolha}
  \tcbitem {\small\color{graymid} Duração}\\ {\Large\bfseries\color{winedark} @@DURATION_HOURS@@h}\\ {\tiny\color{graymid} @@DURATION_MINUTES@@ minutos}
  \tcbitem {\small\color{graymid} Alternativas}\\ {\Large\bfseries\color{winedark} A - D}\\ {\tiny\color{graymid} 4 opções por questão}
\end{tcbitemize}

\vspace{20pt}

% ── "COMO FUNCIONA" com 4 passos e círculos numerados ──
\textbf{\textcolor{winedark}{COMO FUNCIONA}}\\[-4pt]
{\color{winemid}\rule{4cm}{2pt}}\\[8pt]
\begin{itemize}[leftmargin=0pt, itemsep=8pt, label={}]
\item \tikz[baseline=-0.6ex]\node[circle, fill=winedark, text=white, font=\bfseries, inner sep=0pt, minimum size=20pt] {1};~~Imprima esta prova e resolva no papel
\item \tikz[baseline=-0.6ex]\node[circle, fill=winedark, text=white, font=\bfseries, inner sep=0pt, minimum size=20pt] {2};~~Volte à plataforma SanarFlix PRO
\item \tikz[baseline=-0.6ex]\node[circle, fill=winedark, text=white, font=\bfseries, inner sep=0pt, minimum size=20pt] {3};~~Preencha o gabarito digital na plataforma
\item \tikz[baseline=-0.6ex]\node[circle, fill=winedark, text=white, font=\bfseries, inner sep=0pt, minimum size=20pt] {4};~~Envie dentro do tempo para entrar no ranking
\end{itemize}

\vspace{6pt}

% ── Caixa "REGRAS IMPORTANTES" ──
\begin{tcolorbox}[colback=rulesboxtint, colframe=winelight, rounded corners, boxrule=1pt]
  \textbf{\textcolor{winedark}{REGRAS IMPORTANTES}}\\[6pt]
  {\color{winemid}$\cdot$}~O tempo de prova começa no momento do download.\\
  {\color{winemid}$\cdot$}~Envie o gabarito dentro do prazo para participar do ranking.\\
  {\color{winemid}$\cdot$}~Questões não respondidas serão registradas como em branco.
\end{tcolorbox}

\vfill

% ── Footer próprio da capa (replica drawFooter(cover, "Capa"), legacyPdfLib.ts:413) ──
\noindent\textcolor{graylight}{\rule{\textwidth}{0.5pt}}\par\vspace{4pt}
\noindent
\begin{minipage}[t]{0.75\textwidth}
  {\fontsize{7}{9}\selectfont\color{graymid} @@SIMULADO_TITLE@@ · SanarFlix PRO · Modo Offline}
\end{minipage}%
\begin{minipage}[t]{0.24\textwidth}
  \raggedleft{\fontsize{7}{9}\selectfont\color{graymid} Capa}
\end{minipage}

\newpage
 (Task 4) desenha a capa e termina em \newpage.
% Depois vem o cabeçalho/rodapé das páginas de questões (equivalente a
% drawHeader/drawFooter, supabase/functions/generate-exam-pdf/legacyPdfLib.ts:474-494)
% e por fim o corpo de questões, montado pela Task 9 e injetado no marcador
% @@QUESTIONS_BODY@@ (concatenação de N fragmentos no padrão de question.tex,
% Task 5 — sem loop LaTeX, tudo já expandido em texto pelo TS).
%
% Placeholders literais @@...@@ substituídos pelo TS (Task 9). Este arquivo
% permanece com os placeholders literais — não editar valores aqui.

\documentclass[11pt,a4paper]{article}
\usepackage{fontspec}
\setmainfont{PlusJakartaSans-Regular.ttf}[
  Path = /opt/fonts/PlusJakartaSans/,
  BoldFont = PlusJakartaSans-Bold.ttf
]
\newfontfamily\jakartamedium{PlusJakartaSans-Medium.ttf}[Path = /opt/fonts/PlusJakartaSans/]
\usepackage{polyglossia}
\setmainlanguage{portuguese}
\usepackage{ucharclasses}
\newfontfamily\greekfallback{NotoSans-Regular.ttf}[Path = /opt/fonts/NotoSans/]
\setTransitionsForGreek{\begingroup\greekfallback}{\endgroup}
\usepackage{microtype}
\usepackage{graphicx}
\usepackage[margin=2cm]{geometry}
\usepackage{fancyhdr}
\usepackage{xcolor}
\usepackage{enumitem}
\usepackage{needspace}
\usepackage[most]{tcolorbox}
\usepackage{tikz}

% Cores exatas equivalentes ao motor pdf-lib atual (supabase/functions/generate-exam-pdf/index.ts:29-36)
\definecolor{winedark}{HTML}{591230}
\definecolor{winemid}{HTML}{6B1C3D}
\definecolor{winelight}{HTML}{8C2E4C}
\definecolor{graylight}{HTML}{EBEBEB}
\definecolor{graymid}{HTML}{8C8C8C}
\definecolor{graybg}{HTML}{F5F5F5}
\definecolor{titlesubtitle}{HTML}{FFD9E6}
\definecolor{titlesubtitle2}{HTML}{FFB2C7}
\definecolor{rulesboxtint}{HTML}{FAF0F2}

\pagestyle{fancy}
\fancyhf{}
% cabeçalho/rodapé completos entram junto com exam.tex (Task 5) — aqui só reset de estilo

% Exceções de hifenização para termos médicos compostos que o pacote de
% PT-BR erra sozinho. Um item por entrada, sílabas separadas por hífen nos
% pontos de quebra válidos. Vazio por padrão — popular sob demanda conforme
% casos ruins aparecerem em revisão editorial.
% Exemplo (comentado): \hyphenation{ele-tro-car-di-o-gra-ma}
\hyphenation{}


\begin{document}

% cover.tex — Capa (recriação visual de drawCoverPage, supabase/functions/generate-exam-pdf/legacyPdfLib.ts:200-414)
% Inputted from exam.tex (Task 5) logo após \begin{document}. Placeholders @@...@@
% são substituídos por dados reais na Task 9 (Node.js). NÃO editar valores aqui —
% este arquivo permanece com os placeholders literais.

% ── Bloco de título com fundo wine ──
\begin{tcolorbox}[colback=winedark, colframe=winemid, sharp corners, boxrule=1pt, width=\textwidth, left=8pt, right=8pt, top=10pt, bottom=10pt]
  {\fontsize{26}{30}\selectfont\bfseries\color{white} PROVA OFFLINE}\\[4pt]
  {\fontsize{14}{16}\selectfont\color{titlesubtitle} @@SIMULADO_TITLE@@}\\[2pt]
  {\fontsize{10}{12}\selectfont\color{titlesubtitle2} Modo Offline · Gabarito Digital}
\end{tcolorbox}

\vspace{18pt}

% ── 3 info-cards lado a lado: Questões / Duração / Alternativas ──
% Nota: \tcbitem só é definido pelo ambiente tcbitemize (que envolve tcbraster
% internamente) — não por tcbraster diretamente. Mesmas opções do brief, apenas
% o nome do ambiente corrigido.
\begin{tcbitemize}[raster columns=3, raster equal height, colback=graybg, colframe=graylight, rounded corners, boxrule=0.5pt]
  \tcbitem {\small\color{graymid} Questões}\\ {\Large\bfseries\color{winedark} @@QUESTIONS_COUNT@@}\\ {\tiny\color{graymid} Múltipla escolha}
  \tcbitem {\small\color{graymid} Duração}\\ {\Large\bfseries\color{winedark} @@DURATION_HOURS@@h}\\ {\tiny\color{graymid} @@DURATION_MINUTES@@ minutos}
  \tcbitem {\small\color{graymid} Alternativas}\\ {\Large\bfseries\color{winedark} A - D}\\ {\tiny\color{graymid} 4 opções por questão}
\end{tcbitemize}

\vspace{20pt}

% ── "COMO FUNCIONA" com 4 passos e círculos numerados ──
\textbf{\textcolor{winedark}{COMO FUNCIONA}}\\[-4pt]
{\color{winemid}\rule{4cm}{2pt}}\\[8pt]
\begin{itemize}[leftmargin=0pt, itemsep=8pt, label={}]
\item \tikz[baseline=-0.6ex]\node[circle, fill=winedark, text=white, font=\bfseries, inner sep=0pt, minimum size=20pt] {1};~~Imprima esta prova e resolva no papel
\item \tikz[baseline=-0.6ex]\node[circle, fill=winedark, text=white, font=\bfseries, inner sep=0pt, minimum size=20pt] {2};~~Volte à plataforma SanarFlix PRO
\item \tikz[baseline=-0.6ex]\node[circle, fill=winedark, text=white, font=\bfseries, inner sep=0pt, minimum size=20pt] {3};~~Preencha o gabarito digital na plataforma
\item \tikz[baseline=-0.6ex]\node[circle, fill=winedark, text=white, font=\bfseries, inner sep=0pt, minimum size=20pt] {4};~~Envie dentro do tempo para entrar no ranking
\end{itemize}

\vspace{6pt}

% ── Caixa "REGRAS IMPORTANTES" ──
\begin{tcolorbox}[colback=rulesboxtint, colframe=winelight, rounded corners, boxrule=1pt]
  \textbf{\textcolor{winedark}{REGRAS IMPORTANTES}}\\[6pt]
  {\color{winemid}$\cdot$}~O tempo de prova começa no momento do download.\\
  {\color{winemid}$\cdot$}~Envie o gabarito dentro do prazo para participar do ranking.\\
  {\color{winemid}$\cdot$}~Questões não respondidas serão registradas como em branco.
\end{tcolorbox}

\vfill

% ── Footer próprio da capa (replica drawFooter(cover, "Capa"), legacyPdfLib.ts:413) ──
\noindent\textcolor{graylight}{\rule{\textwidth}{0.5pt}}\par\vspace{4pt}
\noindent
\begin{minipage}[t]{0.75\textwidth}
  {\fontsize{7}{9}\selectfont\color{graymid} @@SIMULADO_TITLE@@ · SanarFlix PRO · Modo Offline}
\end{minipage}%
\begin{minipage}[t]{0.24\textwidth}
  \raggedleft{\fontsize{7}{9}\selectfont\color{graymid} Capa}
\end{minipage}

\newpage


% ── Cabeçalho/rodapé para as páginas de questões (equivalente a
% drawHeader/drawFooter, legacyPdfLib.ts:474-494) ──
% Nota: o motor antigo desenha uma barra wine-dark de fundo atrás do cabeçalho
% (drawHeader, legacyPdfLib.ts:474-486); fancyhdr por si só só posiciona texto,
% não pinta um retângulo de fundo. Reproduzir a barra colorida exigiria um
% mecanismo adicional (ex. \fancypagestyle customizado com tikz overlay) fora
% do escopo deste brief — mantido como texto simples por ora; se a aprovação
% visual (gate desta task) pedir a barra de fundo, isso vira um ajuste de
% rodada de fix.
\fancyhead[L]{\small\bfseries sanarflix PRO \quad {\scriptsize S I M U L A D O S}}
\fancyhead[R]{\small @@EXAM_LABEL@@}
\fancyfoot[L]{\scriptsize @@FOOTER_LABEL@@}
\fancyfoot[R]{\scriptsize \thepage}
\renewcommand{\headrulewidth}{0.5pt}
\renewcommand{\footrulewidth}{0.5pt}

@@QUESTIONS_BODY@@

\end{document}
